\section{案例：量化系统}
\subsection{构建}
\begin{frame}{提纲}
\begin{itemize}
\item 构建：插件、角色、检查关卡、案例运行与生产边界。
\item 演进：记录失败、比较、批准、回滚与监测。
\end{itemize}
\end{frame}
\begin{frame}{插件结构与入口}
\begin{center}
\lstinputlisting[style=filetext,basicstyle=\ttfamily\footnotesize,breakatwhitespace=false]{figures/qt_layout.txt}

\end{center}
\begin{itemize}
\item 用 \texttt{/qt:start} 开展研究，用 \texttt{/qt:evo} 改进系统。
\item 用 \texttt{/plugins install <qt>} 安装。修改源码后，重新安装并运行 \texttt{/reload}。
\item 将安装后的代码与工作区 \texttt{.qt/} 中的案例状态、提示词版本分开保存。
\end{itemize}
\end{frame}
\begin{frame}{负责人与专项 Agent}
\begin{center}
\lstinputlisting[style=filetext,basicstyle=\ttfamily\footnotesize,breakatwhitespace=false]{figures/qt_agent.txt}

\end{center}
\begin{itemize}
\item 在角色卡中定义负责人的工具权限及允许调用的专项 Agent。
\item 实际运行时，检查宿主的实际分派与响应记录。
\item 仅凭报告文件，无法证明各专项 Agent 曾在独立上下文中运行。
\end{itemize}
\end{frame}
\begin{frame}{Orchestrator 与检查关卡}
\begin{itemize}
\item Orchestrator 绑定请求，核验 Generator 报告与数值关卡，分派新 Critic 审查并拥有状态转换。
\item 前进前必须有三份匹配的 Generator 报告和三份新的 Critic 审查。
\item 运行 \texttt{python3 <qt>/scripts/check.py} 检查控制器规则。
\end{itemize}
\begingroup
\setstretch{1.0}
\lstinputlisting[style=filetext,basicstyle=\ttfamily\footnotesize]{figures/ca_controller_code.txt}
\endgroup
\par
\end{frame}
\begin{frame}{Hook}
\begin{center}
\lstinputlisting[style=filetext,basicstyle=\ttfamily\footnotesize]{figures/ca_hook_manifest.txt}

\end{center}
\begin{itemize}
\item 可选 Hook 练习（未随插件提供）：将 \texttt{audit\_event.py} 复制到 \texttt{<qt>/hooks/}，合并清单配置，再重新安装并加载。
\item 事件名、工具名及会话与调用 ID 只记录一次；它既不评估候选策略，也不提供完整的专项 Agent 执行轨迹。
\item Kimi 记录 Hook 错误或超时后会继续运行。必要检查仍由控制器执行。
\end{itemize}
\end{frame}
\begin{frame}{运行一个研究案例}
\begin{itemize}
\item 跟踪 \texttt{request\_id}、报告、数值输出、路由、\texttt{trace.jsonl} 和终止结果。
\item 按决策日期选择当时可用的时点版本；重述版本在披露前不可用。
\item 动量契约使用21个有效收盘价：$s_{i,t}=P_{i,t}/P_{i,t-20}-1$，示例为 $12/10-1=20\%$。
\item 在案例前检查稳定排序与最早可执行时间。
\item 样例只测试动量候选方案；仅\texttt{lookback} 和 \texttt{top\_n} 变化。
\item 脚本演练需明确标注；声称真实团队运行时必须出示宿主分派记录。
\end{itemize}
\end{frame}
\begin{frame}{修复被退回的提交\handsup}
\begin{itemize}
\item 缺少专项报告：否决提交，留在当前角色并消耗一次尝试。
\item 请求 ID 过期：执行同样的否决与重试规则；第二次无效提交变为 \texttt{BLOCKED}。
\item 报告完整匹配后，先检查证据再前进。
\item 问题：在临时测试案例中制造缺少报告和过期 ID，哪些修复才能继续？
\item 不带 \texttt{--response} 调用 \texttt{advance}；重新读取 \texttt{request.json}，让所有报告匹配新 ID。检查哈希并保留封存证据。
\end{itemize}
\end{frame}
\begin{frame}{重放、审计与恢复\handsup}
\begin{itemize}
\item 重建：检查输入/代码、已解决与未解决问题项、停止规则、角色卡/版本以及批准记录。
\item 复算：固定输入、代码和配置，比较数值结果。
\item 新的模型运行：使用新案例，保留自己的决定、工具路径和输出。
\item 恢复或复用：核验请求 ID 与哈希；恢复待办工作；已完成案例返回保存结果。
\item 问题：重建本案例。案例文件或单独批准记录是否标明了批准该版本的人？
\end{itemize}
\end{frame}
\begin{frame}{生产数据}
\begin{itemize}
\item 输入：有许可证、带日期的市场与财务数据。
\item 股票池：历史成分、停牌与退市。
\item 执行：相关板块和日期的手数、价格限制、T+1规则与费用。
\item 运营：宿主隔离、责任人，以及独立测试的适配器。
\item 课堂样例使用五只虚构股票和合成工作日数据；生产数据接入、公司行为、组合控制和券商适配器仍待构建。
\end{itemize}
\end{frame}
\begin{frame}{研究验证}
\begin{center}
\resizebox{.94\textwidth}{!}{\begin{tikzpicture}[
  box/.style={draw=blue,fill=blue!5,rounded corners=2pt,text width=3.3cm,minimum height=1.5cm,align=center,font=\small,inner sep=5pt},
  link/.style={-{Stealth[length=5pt]},thick,draw=black!60}]
\node[box] (n0) at (1.7,0) {固定 $t$ 日及之前数据};
\node[box] (n1) at (6.6,0) {修改 $t$ 日之后价格};
\draw[link] (n0) -- (n1);
\node[box,draw=forestgreen,fill=forestgreen!6] (n2) at (11.5,0) {此前的信号\\必须保持不变};
\draw[link] (n1) -- (n2);
\end{tikzpicture}
}
\end{center}
\begin{itemize}
\item 检验未来价格变动、数据重述、更高成本，以及依次推进的样本外区间表现。
\item 通过这些检查，并不能证明策略将获得超额收益。
\end{itemize}
\end{frame}
\begin{frame}{资金授权}
\begin{center}
{\small\textbf{生产环境的目标流程}}\par\smallskip
\resizebox{.94\textwidth}{!}{\begin{tikzpicture}[b/.style={draw=blue,fill=blue!6,rounded corners,text width=3.4cm,minimum height=1.0cm,align=center,font=\small},a/.style={-{Stealth},thick}]
\node[b] (n0) {\textbf{DRAFT}\\带版本号的建议单};
\node[b,right=0.5cm of n0] (n1) {\textbf{REVIEWED}\\人工复核与审批};
\draw[a] (n0)--(n1);
\node[b,right=0.5cm of n1] (n2) {\textbf{RECONCILED}\\导入实际成交对账};
\draw[a] (n1)--(n2);
\end{tikzpicture}
}
\end{center}
\begin{itemize}
\item 检查组合限额、当前持仓、批准有效期与实际成交。
\item 由指定负责人批准资金使用；研究验收通过不等于获得下单授权。
\item \texttt{trade\_sheet.csv} 使用独立的历史合成仓位计算样例，并非来自已验收的候选策略。QT 没有券商连接。
\end{itemize}
\end{frame}
\begin{frame}{记录一次失败}
\begin{itemize}
\item 回到研究工作流，\texttt{/qt:evo} 开始系统改进；\texttt{evo-prepare} 关联完成案例与变更前提示词版本。
\item 记录失败、所需行为、开发案例和独立的最终评估案例。
\item GEPA 等外部候选工具可以提出改动，但未集成到 QT。
\end{itemize}
\lstinputlisting[style=shell,basicstyle=\ttfamily\footnotesize]{figures/ca_capture_command.txt}
\par
\end{frame}
\subsection{演进}
\begin{frame}{提示词改动\handsup}
\begin{center}
\lstinputlisting[style=diff,basicstyle=\ttfamily\small]{figures/timing_prompt_diff.txt}

\end{center}
\begin{itemize}
\item 问题：哪条指令针对了时序遗漏？需要什么证据才能说明它有效？
\item 该 QT 接口只修改 Generator 提示词。策略参数、风险限额、验收规则与控制器代码均不在其修改范围内。
\end{itemize}
\end{frame}
\begin{frame}{比较输出\handsup}
\begin{center}
\small\begin{tabular}{@{}P{4.0cm}P{8.5cm}@{}}
\toprule
\textbf{准备好的对照评测} & \textbf{核验内容} \\
\midrule
原提示词与候选提示词 & 案例与模型设置相同 \\
\rowcolor{gray!10}
Python 检查函数 & 时间顺序与证据引用 \\
输出对照表 & 逐项失败原因与原始输出 \\
\bottomrule
\end{tabular}

\end{center}
\begin{itemize}
\item 交易从 30 Sep 09:30 (UTC+8) 起才允许执行。仅留存了 \texttt{PRICE\_01}。
\item 运行配套检查，或在 promptfoo 中加载 \texttt{promptfooconfig.json}，比较两份输出。
\item 问题：是否有任一输出同时通过两项检查？
\end{itemize}
\end{frame}
\begin{frame}{判断证据\handsup}
\begin{center}
\small\begin{tabular}{@{}P{3.7cm}P{4.3cm}P{4.3cm}@{}}
\toprule
\textbf{检查} & \textbf{原版本} & \textbf{候选版本} \\
\midrule
时点 & 未通过 & 通过 \\
\rowcolor{gray!10}
证据引用 & 通过 & 未通过 \\
采纳判断 & 否决 & 否决 \\
\bottomrule
\end{tabular}

\end{center}
\begin{itemize}
\item 候选版本修正了时序，却引用了缺失的来源 \texttt{MISSING\_02}。
\item 问题：采用该版本前，还需修复哪些问题，并进行哪些独立检查？
\end{itemize}
\end{frame}
\begin{frame}{版本控制}
\begin{itemize}
\item Git 记录源代码历史以及产生候选改动的差异。
\item QT 注册表记录提示词字节、活动指针、父版本和运行配置。
\item MLflow 是外部对应工具，用于提示词版本、评估与追踪；未集成到 QT。
\item 别名用于选择版本，但批准需要独立记录。
\end{itemize}
\end{frame}
\begin{frame}{评估、测试、批准\handsup}
\begin{itemize}
\item (1) 评估：使用提案、审查和回归结果运行 \texttt{evo-evaluate}。
\item (2) 测试：检查代码回归和行为比较，包括留出的案例。
\item (3) 批准：只有明确的人类批准才能运行 \texttt{evo-adopt}；否决则保留活动版本。
\item 研究验收、版本采用和资金授权仍是不同决定。
\end{itemize}
\end{frame}
\begin{frame}{回滚与监测}
\begin{center}
\resizebox{.94\textwidth}{!}{\begin{tikzpicture}[
  box/.style={draw=blue,fill=blue!5,rounded corners=2pt,text width=3.3cm,minimum height=1.5cm,align=center,font=\small,inner sep=5pt},
  link/.style={-{Stealth[length=5pt]},thick,draw=black!60}]
\node[box] (n0) at (1.7,0) {\textbf{批准回滚}\\保留新旧版本记录};
\node[box] (n1) at (6.6,0) {\textbf{恢复版本指针}\\切回先前提示词};
\draw[link] (n0) -- (n1);
\node[box,draw=forestgreen,fill=forestgreen!6] (n2) at (11.5,0) {\textbf{新建案例验证}\\确认行为恢复正常};
\draw[link] (n1) -- (n2);
\end{tikzpicture}
}
\end{center}
\begin{itemize}
\item 获得批准后，用 \texttt{evo-rollback} 恢复此前的指针。保留新旧两个版本及其证据。
\item 已有案例保留原配置，运行中途不切换版本。
\item 在新案例中核验恢复后的版本，再监测后续运行中是否反复出现失败。
\end{itemize}
\end{frame}
\begin{frame}{最终审查\handsup}
\begin{itemize}
\item 证据：追溯并复算已记录工作。
\item 控制：要求失败时阻断或解释。
\item 连续性：恢复已核验工作，避免重复不确定操作。
\item 改进：评估并恢复受治理的系统版本。
\item 选择一项控制，说明它回应谁的需要，并说明还需构建什么。
\end{itemize}
\end{frame}
